\section{The native quadratic diagonal has nonnilpotent tail action}
\label{sec:136-quadratic-tail-action}

\paragraph{Archive, inputs, and exact scope.}
The uniform hypotheses are cote136, physical pp.~117 and119,
footers116 and118, scans117 and119.  The common-source question
is physical p.~121, footer120, scan121.  This module answers a
specific remaining question in the coherent native repair: the
odd quadratic homology need not be annihilated by any power of
a coefficient tail.  It also records a positive refinement of
the graph-residue criterion.  Neither result asserts a nonzero
actual graph residue or resolves integral connectedness.

The inputs are the internally proved binary-digit resolution in
ARCH-136-ACTUAL-TAIL-SYZYGY, the actual native block and graph
cone convention in ARCH-136-NATIVE-MIXED-TOR/RESIDUE, and the
ordinary coefficient comparison used in
ARCH-136-SKELETON-TATE-LIFTS.  The exact norm sequences and all
connecting maps used below are proved here.  No outside preprint
or additional published theorem is invoked.  The already checked
finite-weight Nakayama input in the resolution retains its
published/background Stacks locator in the source module.
The norm argument was developed jointly with the parent; a
separate exact finite calculation by reviewer123 is corroboration,
not a premise of this proof.

\subsection{The free quadratic norm sequences}

Let \(S=\Gamma_{\mathbb Z}(T)\), let \(M=2^a\), and put
\[
 I=I_M=S_{\geq M},\qquad J=(2,I),\qquad
 F=S/J,\qquad \bar S=S/2S.
                                                               \tag{136.QT1}
\]
The same argument works after localization at two.  Write
\(x_i=T^{[M2^i]}\).  There are odd integers \(\kappa_i\)
such that
\[
 x_i^2=2\kappa_i x_{i+1}.
                                                               \tag{136.QT2}
\]
Consider the homological special dg-PD algebra
\[
 P=S\otimes\Lambda(\theta,e_i:i\geq0)
           \otimes\Gamma(f_i:i\geq0),
 \quad
 \partial\theta=2,\quad \partial e_i=x_i,\quad
 \partial f_i=x_i e_i-2\kappa_i e_{i+1}.
                                                               \tag{136.QT3}
\]
Here \(\theta,e_i\) have homological degree one, \(f_i\)
has degree two, and their coefficient weights are respectively
\(0,M2^i,M2^{i+1}\).  Its augmentation is a free resolution
of \(F\).  Indeed, at the prime two this is exactly the checked
binary resolution followed by the two-term resolution for two.
At every odd prime, multiplication by \(\theta/2\) contracts
it.  Its degree-zero cokernel is \(S/(2,x_i:i\geq0)=F\),
since the binary digits of weight at least \(M\) generate
\(I\) modulo two.  In each weight the complex is finite free
and bounded; localization at all primes therefore proves the
assertion integrally as well.

For a free homological complex \(P\), put
\[
 D=P\otimes_S P,\qquad
 \tau(u\otimes v)=(-1)^{|u||v|}v\otimes u,
 \quad G=\ker(1-\tau),\quad L=\ker(1+\tau).
\]
The coefficients and all terms of \(P\) are two-torsion-free.
Thus \(G\) is naturally the arity-two special graded divided
power complex \(\Gamma^2_{\mathrm{gr}}(P)\): its diagonal
basis element for an even generator \(u\) is represented by
\(u\otimes u\), and its off-diagonal basis elements by graded
symmetric sums.  Here \(\gamma_2(1)\) denotes the divided
square of the formal basis vector of \(P_0\), not a divided
power of the coefficient unit.  The complex \(L\) has graded alternating sums
and one diagonal basis element for each odd generator.  Defining
these two complexes inside \(D\) fixes all differential signs.
Set
\[
 A=\operatorname{im}(1+\tau),\qquad
 B=\operatorname{im}(1-\tau),\qquad
 C_e=G/A,\qquad C_o=L/B.
\]
There are four short exact sequences of complexes:
\[
 \begin{gathered}
 0\longrightarrow L\longrightarrow D\longrightarrow A\longrightarrow0,\\
 0\longrightarrow A\longrightarrow G\longrightarrow C_e\longrightarrow0,\\
 0\longrightarrow G\longrightarrow D\longrightarrow B\longrightarrow0,\\
 0\longrightarrow B\longrightarrow L\longrightarrow C_o\longrightarrow0.
 \end{gathered}                                                    \tag{136.QT4}
\]
These follow basiswise: the norms are isomorphisms onto the
appropriate off-diagonal summands and multiply the relevant
diagonal summands by two.  The quotient complexes have zero
differential.  More precisely, \(C_e\) has the Frobenius base
change of \(P_{2j}/2\) in degree \(4j\), and \(C_o\) has
the Frobenius base change of \(P_{2j+1}/2\) in degree \(4j+2\).
For example
\[
 (C_o)_2=\bar S\,\bar\theta\ \oplus
                      \bigoplus_{i\geq0}\bar S\,\bar e_i.
                                                               \tag{136.QT5}
\]
A basis change with coefficient \(c\) acts on these diagonal
quotients with coefficient \(c^2\); this explains the Frobenius
qualification.  For the displayed fixed basis, they are the
indicated free \(\bar S\)-modules; each diagonal label has
twice the weight of its original basis vector.  There is no claim that
an arbitrary coefficient-linear differential survives unchanged
under that base change.

Finally
\[
 H_j(D)=\operatorname{Tor}_j^S(F,F),\qquad
                       J H_j(D)=0\quad(j\geq0).
                                                               \tag{136.QT6}
\]
The first assertion follows by tensoring two free resolutions.
For the second, multiplication by an element of \(J\) on a
free resolution lifts the zero map of \(F\), and is therefore
chain homotopic to zero; tensoring that homotopy proves the
annihilation.  This also supplies an elementary proof of the
specific Tor action used below.

\subsection{The odd norm obstruction}

\paragraph{ARCH-136-QUADRATIC-TAIL-ACTION.}
For the resolution (136.QT3), there is a coefficient-linear,
weight-preserving map
\[
 \Phi_5:H_5\bigl(\Gamma^2_{\mathrm{gr}}(P)\bigr)
                                  \longrightarrow(C_o)_2
 \quad\text{such that}\quad
                 I^4(C_o)_2\subseteq\operatorname{im}\Phi_5.
                                                               \tag{136.QT7}
\]
Moreover
\[
              4H_j\bigl(\Gamma^2_{\mathrm{gr}}(P)\bigr)=0
                                                   \quad(j>0).
                                                               \tag{136.QT8}
\]
In particular no power of any coefficient tail annihilates
\(H_5(\Gamma^2_{\mathrm{gr}}(P))\).  There is a class whose
\(\bar\theta\)-coordinate under \(\Phi_5\) is
\(T^{[15M]}\) modulo two.  Multiplication of that class by
every sufficiently high fresh binary digit remains nonzero.

\begin{proof}
First construct a map in degree three.  The connecting map
\[
                    \delta_o:(C_o)_2\longrightarrow H_1(B)
\]
from the fourth sequence in (136.QT4), followed by the connecting
map \(H_1(B)\to H_0(G)\) from the third, is particularly simple.
Let \(y\in P_1\) be a basis element and write
\(\partial y=c\cdot1\).  Represent its diagonal class in
\(L_2\) by \(y\otimes y\).  Its boundary is
\(c(1\otimes y-y\otimes1)\).  A norm preimage in \(D_1\)
is \(c(1\otimes y)\), whose boundary is
\(c^2(1\otimes1)\).  Hence the composite is
\[
       \bar y\longmapsto c^2\gamma_2(1)
                   \quad\text{in}\quad
       H_0(G)=S/(4,2x_i:i\geq0)=S/(4,2I).
                                                               \tag{136.QT9}
\]
For \(y=\theta\) its value is four; for \(y=e_i\) it is
\(x_i^2=2\kappa_i x_{i+1}\).  All values are zero in that
quotient.  Exactness of the third sequence thus puts the image
of \(\delta_o\) inside the image of \(H_1(D)\), which is
annihilated by \(I\) by (136.QT6).  Exactness of the fourth
sequence now gives
\[
 I(C_o)_2\subseteq
     \operatorname{im}\bigl(H_2(L)\longrightarrow(C_o)_2\bigr).
                                                               \tag{136.QT10}
\]
Multiplying a lift in \(H_2(L)\) by one more element of \(I\)
kills its image in \(H_2(D)\).  The first sequence consequently
lifts it to \(H_3(A)\).  Thus the composite
\[
           H_3(A)\longrightarrow H_2(L)\longrightarrow(C_o)_2
\]
has image containing \(I^2(C_o)_2\).

This composite factors through \(H_3(G)\).  Indeed, the second
sequence makes \(H_3(A)\to H_3(G)\) onto, with kernel the
image of \((C_e)_4\).  To compute that kernel's image in
\((C_o)_2\), take an even basis element \(v\in P_2\).
The diagonal \(v\otimes v\) represents its class in
\((C_e)_4\); its boundary has the norm preimage
\(v\otimes\partial v\).  The boundary of that preimage is
\(\partial v\otimes\partial v\).  Projecting to \((C_o)_2\)
therefore gives the squares, modulo two, of the matrix
coefficients of \(\partial:P_2\to P_1\).  All of those
coefficients lie in \((2,x_i:i\geq0)\), and every such
coefficient has square zero modulo two.  The latter assertion
follows directly from (136.QT2) and additivity of squaring in
characteristic two.  The composite on \((C_e)_4\) is zero.
We have proved the existence of a well-defined linear map
\[
        \Phi_3:H_3(G)\longrightarrow(C_o)_2,
          \qquad I^2(C_o)_2\subseteq\operatorname{im}\Phi_3.
                                                               \tag{136.QT11}
\]

The supports of \(C_e\) and \(C_o\) give canonical
isomorphisms
\[
                      H_5(A)\simeq H_5(G),\qquad
                      H_4(B)\simeq H_4(L).
\]
Using the connecting maps from the first and third sequences
in (136.QT4), define
\[
 H_5(G)\simeq H_5(A)\longrightarrow H_4(L)
       \simeq H_4(B)\longrightarrow H_3(G)
                            \xrightarrow{\Phi_3}(C_o)_2.
                                                               \tag{136.QT12}
\]
This is \(\Phi_5\).  Its first four arrows have image
containing \(I^2H_3(G)\): multiply a given class in \(H_3(G)\)
by one tail element to kill its image in \(H_3(D)\), lift
it to \(H_4(B)\), and multiply by a second tail element to
kill the resulting image in \(H_4(D)\); exactness then lifts
it to \(H_5(A)\).  Applying (136.QT11) proves (136.QT7).
These are existence of lifts along exact module maps, not a
choice of a coefficient-linear family of lifts.

Multiplication by \(\theta\) on \(P\) is a homotopy for
\(2\operatorname{id}_P\).  Extend that degree-one linear map
as a special derivation of the free graded PD algebra on the
underlying complex \(P\).  Its commutator with the differential
is twice the arity derivation.  On arity two it is therefore
\(4\operatorname{id}\), proving (136.QT8), including all
normalizations of divided squares.

Finally the product of the four tail digits of weights
\(M,2M,4M,8M\) is \(T^{[15M]}\) modulo two, by the
no-carry binomial formula.  Use its \(\bar\theta\)-coordinate
in (136.QT7) and set the other coordinates to zero.  A lift
has nonzero image after multiplication by any digit of weight
\(2^bM\) with \(b\geq4\), and likewise after products of
arbitrarily many distinct such digits.  Choosing all their
weights beyond any specified tail proves that no power of
that tail annihilates the homology.  The maps preserve weights,
so these are also nonzero classes in unbounded coefficient
weights.  This proves every assertion.
\end{proof}

\paragraph{Native interpretation and its limit.}
Apply the native comparison
ARCH-136-DIAGONAL-REDUCTION first over \(\mathbb Z_{(2)}\),
choosing the binary resolution of \(I\) there.  In that native
block the arity-two complex \(K\) contains \(e=ab\),
\(z=\gamma_2(b)\), and all coherent parameters.  Reversing
its cochain degrees makes it the corresponding resolution of
\(F\) truncated at its fourth syzygy.  The quadratic
truncation comparison identifies its cochain-degree-three
quadratic homology with the homological degree-five calculation
above.  The actual native arity-four multiplication sends
\(\gamma_2(z)\) to \(3\gamma_4(b)\) and \(ez\) to
\(3a\gamma_3(b)\), and is the evident identity on the other
summands.  It is an isomorphism at two; its only nonunit
integral multipliers occur in form degrees seven and eight.
The lower native arities have no unwanted \(H^0\) or \(H^1\).
In the homogeneous native model, the auxiliary-arity-zero
part is just \(S\) in form zero.  Thus its factorwise RH1
changes no form-degree-two-or-higher term in arity four.

For an arbitrary integral free resolution of \(I_M\), do not
identify its native complex directly with (136.QT3).  Localize
that resolution at two.  The binary minimal resolution is its
strict chain retract by the comparison argument in
ARCH-136-ACTUAL-TAIL-SYZYGY.  Applying the linear resolution
maps to every \(U_i,V_i\), and fixing the native \(a,b\),
gives a strict retraction of the complete coherent native
algebras.  Their bottom relations are preserved because the
maps commute with the resolution differential and augmentation.
The local nonannihilated arity-four class therefore remains
nonannihilated in the localized integral native algebra.
The same lower-arity argument for homogeneous RH1 preserves
its positive form terms.  If an integral tail power annihilated
that native cohomology, its localization would do so too,
a contradiction.  This transfer uses an actual strict retract,
not invariance of \(\Gamma\) under arbitrary quasi-isomorphisms.
All native weights receive the common shift \(4m\), where
\(m\) is the original input weight.

Consequently the proposed proof that all these diagonal graph
residues vanish by \(I\)-power annihilation cannot work.  A
particular target's residue may still vanish for another reason.
No map from these surviving classes to a nonzero graph class
has been constructed here, and no obstruction to every possible
continuation or common source is claimed.

\subsection{The torsion contribution to the graph cone}

\paragraph{ARCH-136-TORSION-GRAPH-RESIDUE.}
Let \(B\) be a coefficient-linear cochain subcomplex of one
source simplex, restricting literally to zero on
\(G=\operatorname{sk}_1\Delta[n]\), and let \(f:B\to C_n\)
be induced by a strict map to an admissible target with the
fixed common graph algebra.  Work above the uniform coefficient
comparison and row-extension bounds through degree three, so
that the graph obstruction is defined as in
ARCH-136-NATIVE-MIXED-RESIDUE.  Use the cone convention
\[
 P_N(B)^r=B^r\oplus B^{r+1},\qquad
                  d(x,y)=(dx-Ny,-dy).
                                                               \tag{136.QT13}
\]
The graph obstruction
\[
             H^2(P_N(B))\longrightarrow H^1(G;S/NS)
\]
is zero on the image of the integer-torsion subgroup of
\(H^2(B)\), first mapped to \(H^2(B)/NH^2(B)\) and
then to the cone cohomology.  In particular, if \(H^2(B)\) is
integer torsion, its entire contribution in
\[
 0\longrightarrow H^2(B)/NH^2(B)
  \longrightarrow H^2(P_N(B))
  \longrightarrow H^3(B)[N]\longrightarrow0
                                                               \tag{136.QT14}
\]
has zero graph obstruction.  Only the right-hand quotient then
requires further control.

\begin{proof}
For a closed \(x\in B^2\), choose \(u\in C_n^1\) with
\(du=f(x)\).  It exists by the target's vertical comparison.
Because \(x\) is zero on the graph, \(u|_G\) is a closed
one-form.  Changing \(u\) changes it by the restriction of a
closed target one-form, hence by an exact graph form because
\(H^1(C_n)=0\) in the stated weights.  If \(x\) changes
by a boundary \(db\) in \(B\), the primitive changes by
\(f(b)\), which is zero on the graph.  This defines a
coefficient-linear homomorphism
\[
                 H^2(B)\longrightarrow H^1(G;S).
                                                               \tag{136.QT15}
\]
Linearity follows by adding or multiplying individual primitives;
no linear selection of primitives is required.  The usual graph
coefficient comparison identifies the target with a finite
free sum of copies of \(S\), and in particular it is
integer-torsion-free.  Therefore (136.QT15) is zero on all
integer-torsion classes.

The cone image of \(x\) is represented by \((x,0)\), and
\((u,0)\) is its absolute cone primitive by (136.QT13).
Its graph class is exactly the reduction modulo \(N\) of
(136.QT15), by naturality of the coefficient and cone maps.
Thus it vanishes on the asserted subgroup.  The exact sequence
(136.QT14) follows directly from the same cone convention.
This proves the result for each target individually, with
one uniform comparison range and every modulus \(N\).
\end{proof}

\paragraph{Verification status.}
The author completed a full saved-file reading, checking the
integral resolution, all four norm sequences, both explicit
connecting-map composites, the two successive lifting arguments,
the arity homotopy, fresh-digit action, and the torsion-cone
refinement.  Reviewer1348 completed a separate whole saved-proof reading
and communicated PASS for both numbered statements.  The
parent completed a contributor whole saved-proof reading of
the norm and residue proofs and communicated PASS; the parent
contributed to both arguments, so this is not independent
rediscovery.  The explicit strict-retract transfer paragraph
was added after those readings and has been reread by the
author and parent; the parent communicated a targeted PASS.
Reviewer123 subsequently completed a whole saved-proof PASS
including that strict-retract paragraph.  The finite
computations are not used as proof.  The separately owned native
bridge is ARCH-136-DIAGONAL-REDUCTION.
